% Pursuing Stacks (cotes 134-2 to 134-8): archive cote 134-5, site page 50 (PDF page 51), typescript page 237.
% typescript. Draft compared: ps4.tex:2483-2552.
% AI-transcribed by Claude Opus 5.5 (claude-opus-5-5, adaptive thinking, effort high), Message Batches API, batch ps2 (msgbatch_01NVwr1Yuo8qJf5rykeGBXPK), 2026-11-12; not yet checked by eye.
\scanpage{134-5/50}
\mspage{237}

ment of the proposition - all six conditions are equivalent in this case.
If no B is given, but still assuming A and $(M,M_{as})$ aspheric, we can
choose a generating subcategory for $(M,M_{as})$ large enough in order to
contain $i(A)$ and $e_M$ , and we get that conditions (i) to (ii') are equi-%
valent, qed.

76. It's time to give some examples.

1) Take $M = (Cat)$, $M_{as} =$ set of $\underline{W}$-aspheric objects in (Cat). We
then got an asphericity structure, as we see by taking any weak test
category A in (Cat)\inserted[typed]{(def.2 on page 172)}, and the functor
\[ i_A \ : \ a \longmapsto A_{/a} \ : \ A \longrightarrow (Cat) \quad , \]
which satisfies indeed (i) and (ii) of p.231 above.\margin{//} We may call
\[ (7) \qquad ((Cat),\ (Cat)_{\underline{W}\text{-as}} \overset{\hand{def}}{=} (Cat)_{as} \ ) \]
the ``basic asphericity structure'' , giving rise (by taking the correspon-%
ding set W of ``weak equivalences'') to the ``basic modelizer''\deleted{\unsure{,}} $((Cat), \underline{W})$.
Turning attention \xout{from} towards the former corresponds to a shift in em-%
phasis; whereas previously our\inserted{main} emphasis has \xout{\illegible{1}} been dwelling consistent-%
ly with the notion of ``weak equivalence''\hand{,} namely with \xout{\illegible{1}} giving on a
category M a bunch of \uline{arrows} W , here we are \xout{\illegible{1}} working rather with the
notion of ``aspheric \uline{objects}'' as the basic notion. One hint in this direc-%
tion comes from prop.1 \xout{\illegible{1}} on p.214 , when we saw that for a functor
$f\colon A \longrightarrow B$ between small categories, giving rise to $f^{*}\colon B\hat{} \longrightarrow A\hat{}$ (a map
\margin{?} between asphericity structures, as a matter of fact), asphericity of f
can \xout{be} always\inserted{be} expressed in manifold ways as a property of f relative to
the notion of aspheric \uline{objects} in $A\hat{}$and $B\hat{}$, but not as a property relati-%
ve to the notion of weak equivalences, unless A and B are assumed to be
aspheric.

Coming back to the case of the\over{x}{``}basic asphericity structure'' (7),
we get more general types of ``aspheric'' functors $A \longrightarrow (Cat)$ than functors
$i_A$ , by taking any weak test category A and any weak test functor
$i \ : \ A \longrightarrow (Cat)$ (cf def.5 p. 175 ),\inserted[typed]{provided however the objects i(a) are aspheric .} It would seem though that,\xout{\illegible{1}}
for a given A, even assuming A to be a \uline{strict} test category say (\over{or}{\unsure{and}} even
a ``contractor''\inserted{moreover}), \xout{\illegible{3}} and restricting to functors
$i\colon A \longrightarrow (Cat)$ which factor through $(Cat)_{as}$ (to make them eligible for
being ``a\unsure{pp|sp}heric functors''), the \deleted{property}\inserted{condition} for i to be a weak test functor\hand{,}
namely for $i^{*}$ to be ``model preserving''\hand{,} is substantially stronger than
mere ``asphericity'': indeed the latter just means that $i^{*}$ transforms weak

\margin{50}
